% Part: sets-functions-relations
% Chapter: functions
% Section: isomorphisms

\documentclass[../../../include/open-logic-section]{subfiles}


\begin{document}

\olfileid[tr]{sfr}{fun}{iso}
\olsection{İzomorfizm}

\begin{explain}
Bir \emph{izomorfizm}, ilişkilendirdiği kümelerin yapısını koruyan bir
bire bir örten fonksiyondur; burada yapı, kümelerin !!{element}s{}ı arasında
geçerli olan bağıntılardan oluşur. $X=\{1,2,3\}$ ve $Y=\{4,5,6\}$ kümelerini
ele alın. Bu iki küme ardıl, küçüktür ve büyüktür bağıntılarıyla
yapılandırılmıştır. İki küme arasındaki bir izomorfizm, bu yapıları koruyan
!!a{bijection}{}dur. Dolayısıyla !!a{bijective} $f \colon X \to Y$ fonksiyonu,
$i<j$ ancak ve ancak $f(i)<f(j)$, $i>j$ ancak ve ancak $f(i)>f(j)$ ve $j$,
$i$'nin ardılı ancak ve ancak $f(j)$, $f(i)$'nin ardılıysa bir izomorfizmdir.
\end{explain}

\begin{defn}[İzomorfizm]
$X$ ile $Y$ kümeler, $R$ ile $S$ de sırasıyla $X$ ve $Y$ üzerinde bağıntılar
olmak üzere $U$ çifti $\langle X, R\rangle$, $V$ çifti ise $\langle Y,
S\rangle$ olsun. $X$'ten $Y$'ye bir !!{bijection} $f$, ancak ve ancak
bağıntısal yapıyı koruyorsa $U$'dan $V$'ye bir \emph{izomorfizm}dir; yani
$X$'teki her $x_{1}$ ve $x_{2}$ için $\tuple{x_1,x_2} \in R$ ancak ve ancak
$\tuple{f(x_1),f(x_2)} \in S$ olur.
\end{defn}

\begin{ex}
$X=\{1,2,3\}$ ile $Y=\{4,5,6\}$ kümelerini ve küçüktür ile büyüktür
bağıntılarını ele alın. $f(x) = 7-x$ ile verilen $f\colon X \to Y$ fonksiyonu,
$\tuple{X,<}$ ile $\tuple{Y,>}$ arasında bir izomorfizmdir.
\end{ex}

\end{document}
